%\documentclass[11pt,a4paper,uplatex,dvipdfmx]{ujarticle} 		% for uplatex
\documentclass[11pt,a4j,dvipdfmx]{jarticle} 					% for platex
\input{pieces/form00_header} % pieces
\input{pieces/kakenhi7} % pieces
\input{pieces/form01_header} % pieces
% ===== Global year-dependent definitions for the Kakenhi form ===========
%　基本情報
\newcommand{\研究開始年度}{2027}
\newcommand{\研究開始元号年度}{09}	%令和

\newcommand{\一年目西暦}{2027}
\newcommand{\二年目西暦}{2028}
\newcommand{\三年目西暦}{2029}
\newcommand{\四年目西暦}{2030}
\newcommand{\五年目西暦}{2031}
\newcommand{\六年目西暦}{2032}

\newcommand{\一年目}{9}
\newcommand{\二年目}{10}
\newcommand{\三年目}{11}
\newcommand{\四年目}{12}
\newcommand{\五年目}{13}
\newcommand{\六年目}{14}

\newcommand{\一年目J}{９}
\newcommand{\二年目J}{１０}
\newcommand{\三年目J}{１１}
\newcommand{\四年目J}{１２}
\newcommand{\五年目J}{１３}
\newcommand{\六年目J}{１４}


 % pieces
\input{pieces/hook3} % pieces
%#Name: kiban_s
\input{pieces/form04_jsps_headers} % pieces
\input{pieces/form04_kiban_s_header} % pieces
% ===== Global definitions for the Kakenhi form ======================
%　基本情報
%
%------ 研究課題名  -------------------------------------------
\newcommand{\研究課題名}{加速器実験による短距離バリオン間相互作用の解明と中性子星ハイペロンパズルの解決}

%----- 研究機関名と研究代表者の氏名-----------------------
\newcommand{\研究機関名}{東北大学}
\newcommand{\研究代表者氏名}{市川裕大}
\newcommand{\me}{\underline{\underline{Y.~Ichikawa}}}
%---- 研究期間の最終年度 ----------------
\newcommand{\研究期間の最終元号年度}{13}  %令和で、半角数字のみ
%========================================

\input{pieces/inst_general_images} % pieces
%inst_kiban_s.tex
\newcommand{\EnvironmentInstructions}{%
	\textcolor{red}{(\texttt{\textbackslash EnvironmentInstructions}をコメントアウトしてください。)}\\
	\includegraphicsFullWidth{subject_headers/inst_env.pdf}
}
 % pieces
% user07_header
% ===== my favorite packages ====================================
% ここに、自分の使いたいパッケージを宣言して下さい。
\usepackage{amsmath}
\usepackage{wrapfig}
\usepackage{capt-of}
\usepackage{amssymb}
\usepackage{tikz}
\usetikzlibrary{positioning,shapes.geometric,arrows.meta,calc}
%\usepackage{mb}
%\DeclareGraphicsRule{.tif}{png}{.png}{`convert #1 `dirname #1`/`basename #1 .tif`.png}
\usepackage{lineno}

% ===== my personal definitions ==================================
% ここに、自分のよく使う記号などを定義して下さい。
\newcommand{\klpionn}{K_L \to \pi^0 \nu \overline{\nu}}
\newcommand{\kppipnn}{K^+ \to \pi^+ \nu \overline{\nu}}

% ----- 業績リスト用 -------------
\newcommand{\paper}[6]{%
	% paper{title}{authors}{journal}{vol}{pages}{year}
	\item ``#1'', #2, #3 {\bf #4}, #5 (#6).			% お好みに合わせて変えてください。
}

\newcommand{\etal}{\textit{et al.\ }}
\newcommand{\ca}[1]{*#1}	% corresponding author;   \ca{\yukawa}  みたいにして使う
\newcommand{\invitedtalk}{招待講演}

\newcommand{\yukawa}{H.~Yukawa}					% no underline
%\newcommand{\yukawa}{\underline{\underline{H.~Yukawa}}}	% with 2 underlines
\newcommand{\tomonaga}{S.~Tomonaga}

\newcommand{\prl}{Phys.\ Rev.\ Lett.\ }		% よく使う雑誌も定義すると楽

% ===== 欄外メモ ==================
\newcommand{\memo}[1]{\marginpar{#1}}
%\renewcommand{\memo}[1]{}	% 全てのメモを表示させないようにするには、行頭の"%"を消す

\input{pieces/hook5} % pieces

\begin{document}
\input{pieces/hook7} % pieces
%#Split: 01_purpose_plan  
%#PieceName: p01_purpose_plan
\input{pieces/p01_purpose_plan_00}
\section{研究目的、研究方法など}
%　　　　＜＜最大　6ページ＞＞

%s02_purpose_plan_with_abstract
\noindent
\textbf{（概要）}\\
%begin 研究目的及び研究計画の概要空行付き ＝＝＝＝＝＝＝＝＝＝＝＝＝＝＝＝＝＝＝＝
	本研究の目的は、極限状態の超高密度天体である中性子星の内部構造を決定づける最大の難問「ハイペロンパズル」の解決に向け、ハドロン加速器実験を用いた独創的な新手法により、ハイペロン（$\Lambda$粒子）と核子の間に働く短距離相互作用を実験的に完全解明することである。
	
	理論的には、真空中における$\Lambda N$間の引力獲得メカニズムとして、$\Lambda$粒子が中間状態でストレンジクォークを含む$\Sigma$粒子へと遷移しながら2つの$\pi$中間子を交換する「$2\pi$交換プロセス」が極めて有力視されている。しかし、中性子星深部のような高密度媒質中においては、密度の上昇に伴い核子のパウリ排他律によってこの引力プロセスが強く抑制（パウリブロッキング）され、結果として引力が消滅し、クォークの重なりに起因する短距離からの強い斥力が支配的になるという仮説が提唱されている。
	
	本研究では、この斥力化シナリオの正否に最終決着をつけるため、J-PARC等のハドロンビームラインを拠点として以下の相補的な研究を展開する。(1) $\Sigma N$カスプ分光（J-PARC E90実験）による$\Lambda N$-$\Sigma N$間のチャンネル結合力の高精度決定、(2) $\alpha\Lambda$フェムトスコピー（J-PARC E88実験およびRHIC-STAR実験）を用いた、原子核中で特異的に中心密度が高い$\alpha$粒子内部の極小ソースによる短距離相互作用の直接抽出、および(3) カイラル有効場の理論（Chiral EFT）に基づく先進的なバリオン多体系理論モデルへの統合解析である。もし本研究の検証により、高密度下での斥力化メカニズムのみではハイペロンパズルを説明し得ないという結論に至った場合、状況証拠的に中性子星深部におけるクォーク物質への非閉じ込め相転移の発生が極めて濃厚となる。このように、本研究はバリオン間相互作用の根源的理解を通じて、極限物質観の劇的な変革に迫る挑戦的な試みである。
	
	\vspace*{2zw}
%end 研究目的及び研究計画の概要空行付き ＝＝＝＝＝＝＝＝＝＝＝＝＝＝＝＝＝＝＝＝

\noindent
\rule{\linewidth}{1pt}\\
\noindent
\textbf{（本文）}\\
本研究が挑む相補的な実験群と理論統合による全体像を図\ref{fig:overview}に示す。
%begin 研究目的と研究計画	 ＝＝＝＝＝＝＝＝＝＝＝＝＝＝＝＝＝＝＝＝

\begin{figure}[h]
  \centering
  \begin{tikzpicture}[
    box/.style={draw, rounded corners=4pt, thick, align=center, inner sep=4pt},
    arrow/.style={->, >=latex, thick, draw=gray!80!black},
    lbl/.style={fill=white, inner sep=1.5pt, font=\scriptsize}
  ]
    % 中央柱（E88/STAR）
    \node[box, fill=orange!5, draw=orange!80!black, text width=4.5cm] (e88) {
      \textbf{$\alpha\Lambda$フェムトスコピー}\\
      {\scriptsize (J-PARC E88 / STAR)}\\
      {\scriptsize ［代表：市川 ／ 分担：佐甲・大浦］}\\
      \rule{4.2cm}{0.4pt}\\
      \textcolor{orange!80!black}{\textbf{高密度下での斥力化実証}}\\
      {\scriptsize 中心密度$2\rho_0$の極小ソース活用}\\
      {\scriptsize パウリブロッキングによる引力抑制と短距離斥力芯の顕在化を抽出}
    };
    % 左柱（E90）
    \node[box, fill=blue!5, draw=blue!60!black, text width=4.5cm, left=0.2cm of e88] (e90) {
      \textbf{$\Sigma N$カスプ分光}\\
      {\scriptsize (J-PARC E90)}\\
      {\scriptsize ［代表：市川 ／ 分担：谷田］}\\
      \rule{4.2cm}{0.4pt}\\
      \textcolor{blue!80!black}{\textbf{真空中の引力源泉確定}}\\
      {\scriptsize $\Lambda N$-$\Sigma N$チャンネル結合力決定}\\
      {\scriptsize 真空中で強い引力を生む$2\pi$交換プロセスの定量的寄与を確立}
    };
    % 右柱（HYPS等）
    \node[box, fill=green!5, draw=green!60!black, text width=4.5cm, right=0.2cm of e88] (hyps) {
      \textbf{$\Lambda p$散乱実験}\\
      {\scriptsize (HYPS / E86等)}\\
      {\scriptsize ［研究協力者：三輪 ほか］}\\
      \rule{4.2cm}{0.4pt}\\
      \textcolor{green!80!black}{\textbf{真空中短距離力の基準}}\\
      {\scriptsize 中高エネルギー散乱データ取得}\\
      {\scriptsize 斥力化の比較基準となる真空中の裸の斥力芯サイズを直接測定}
    };

    % タイトルテキスト
    \node[font=\footnotesize\bfseries, above=0.2cm of e88] (exp_title) {独創的な3つの相補的実験アプローチ};

    % 理論研究ノード
    \node[box, fill=brown!5, draw=brown!60!black, text width=12cm, below=1.2cm of e88] (theory) {
      \textbf{理論研究} {\footnotesize （カイラル有効場の理論 Chiral EFT）} {\scriptsize ［分担：神野・神谷］}\\
      {\scriptsize 実験データを統合した解析により、高密度まで適用可能な「密度依存性バリオン間相互作用モデル」を構築}
    };

    % 最終目標ノード
    \node[box, fill=red!5, draw=red!70!black, very thick, text width=13cm, below=1.0cm of theory] (goal) {
      \textbf{最終目標：中性子星ハイペロンパズルの解明}\\
      \rule{12cm}{0.4pt}\\
      {\footnotesize \textcolor{red!70!black}{\textbf{斥力化シナリオの実証}}：高密度でのパウリ抑制によるEoS硬化 $\to$ $2M_\odot$中性子星を支持}\\
      {\scriptsize （※斥力が不十分と判明した場合は、状況証拠的に「クォーク物質への相転移シナリオ」へと決着）}
    };

    % 矢印の描画
    \draw[arrow] (e90.south) -- node[lbl, pos=0.4] {引力・結合力} (theory.north -| e90.south);
    \draw[arrow] (e88.south) -- node[lbl, pos=0.4] {高密度斥力芯} (theory.north);
    \draw[arrow] (hyps.south) -- node[lbl, pos=0.4] {真空中基準力} (theory.north -| hyps.south);
    
    \draw[arrow, very thick, draw=red!70!black] (theory.south) -- node[lbl, text=red!70!black] {極限状態方程式(EoS)の導出} (goal.north);

  \end{tikzpicture}
  \caption{本研究の全体像。真空中での引力メカニズム（$2\pi$交換）と高密度媒質中でのパウリブロッキングによる斥力化を実験的に抽出・比較し、理論統合を通じてハイペロンパズルの解決に挑む。}
  \label{fig:overview}
  \vspace*{1zw}
\end{figure}

\subsection{本研究の学術的背景や着想に至った経緯、研究課題の核心をなす学術的「問い」}
\subsubsection{学術的背景：中性子星とハイペロンパズルの深い陥穽}
私たちの存在するこの宇宙には、通常の物質密度の数倍に達する極限的な超高密度状態が実現している天体が存在する。それが「中性子星」である。中性子星はまさに宇宙に浮かぶ巨大な原子核とも呼ぶべき存在であり、その中心部においては標準原子核密度（$\rho_0 \sim 0.17\,\text{fm}^{-3}$、質量換算で数億トン/\text{cm}$^3$）を遥かに超える極限環境が形成されている。この地上実験室では直接再現することのできない超高密度物質において、物質がどのような形態をとり、どのようなハドロン間相互作用によって巨大な重力崩壊を支えているのかを解明することは、現代の核物理学および天体物理学における最重要課題の一つである。

通常の原子核は、アップ（$u$）クォークとダウン（$d$）クォークからなる陽子と中性子（核子）のみによって構成されている。しかし中性子星内部に向かって密度が高まると、それに伴い核子のフェルミエネルギーが著しく増大する。ある閾値を超えると、エネルギーの高い中性子をさらに系に詰め込むよりも、第三のフレーバーであるストレンジ（$s$）クォークを含む重粒子「ハイペロン」を生成した方が系の全エネルギーが低くなるため、中性子星深部におけるハイペロンの出現は熱力学的・物理的な必然となる。特に、最も軽いハイペロンである$\Lambda$粒子については、これまでの$\Lambda$ハイパー核の系統的な精密分光研究から、標準原子核密度$\rho_0$において$\Lambda$粒子が感じるポテンシャルの深さが約 $-30\,\text{MeV}$ であることが確定しており、$\Lambda$粒子と核子の間には明確な引力が働くことが実証されている。したがって、$\Lambda$粒子は中性子星中心部に極めて出現しやすい状態にある。

しかし、核物理学の標準的な理論モデルに基づく状態方程式（Equation of State; EoS）の予測によれば、ハイペロンの出現はフェルミ圧を減少させ、状態方程式を劇的に軟化（圧力を低下）させる。その結果、現在までに天体観測によって確固たる証拠が得られている太陽質量の2倍（$2M_\odot$）もの重い中性子星の重力を支えることができず、ブラックホールへと崩壊してしまうという深刻な理論的矛盾が生じる。これが広く知られる「ハイペロンパズル」である。我々核物理学を専攻する研究者の立場からは、ミクロな相互作用の素過程に基づき、この天体スケールの矛盾を根本的に解決しなければならない。

\subsubsection{パズル解決に向けた2つの有力な仮説と本研究のアプローチ}
現在、このハイペロンパズルを解決し得るシナリオとして、国際的に以下の2つの有力な仮説が議論されている。
\begin{enumerate}
    \item \textbf{相互作用の斥力化シナリオ}：真空中あるいは標準原子核密度$\rho_0$においては引力的である$\Lambda N$相互作用が、中性子星深部のような高密度環境下においては強い「斥力」へと転化するため、$\Lambda$粒子の出現自体がエネルギー的に抑制される、あるいは出現したとしても多体系全体として強い圧力を生み出す硬いEoSが維持されるという仮説。これを立証するためには、高密度媒質中で$\Lambda N$相互作用が斥力化するメカニズムを定量的かつミクロに解明する必要がある。
    \item \textbf{クォーク物質への非閉じ込め相転移シナリオ}：ハイペロンが出現する密度領域に達する直前にハドロン物質からの相転移が起き、ハドロンの閉じ込めから解放されたクォーク物質（クォーク・グルーオン・プラズマ的物質）が形成されることで、$2M_\odot$以上の質量においても強い圧力を維持するという仮説。
\end{enumerate}

\subsubsection{研究課題の核心をなす学術的「問い」と着想に至った経緯}
本研究が挑む核心的な学術的「問い」は、まさに\textbf{「核物理は中性子星におけるハイペロンパズルを解くことができるのか？」}という根源的な命題にある。

核物理学の観点からは、超高密度下において物質がドラスティックに相転移しクォーク物質となるシナリオは極めて魅力的であり、そのような劇的な現象が実際に星の内部で起きているのかを直接検証したいという強い動機がある。しかしながら、中性子星内部のような「低温かつ原子核密度の数倍に達する高密度環境」を地上実験室で完全に再現することは現状不可能である（高温・高密度状態であれば、中高エネルギー領域の重イオン衝突実験で部分的に再現可能である）。そのため、クォーク相転移の有無を地上の実験から直接証明することは極めて困難である。

そこで申請者は、前者の仮説である\textbf{「高密度環境において$\Lambda$粒子が強い斥力を感じ得るのか？」}という検証可能な疑問にフォーカスする着想に至った。すなわち、加速器ビームを用いた多角的な実験データに基づき、相互作用の斥力化シナリオによってハイペロンパズルを矛盾なく説明できるか否かを徹底検証する。もし本研究による精密な相互作用の決定を通じ、斥力化メカニズムのみでは重い中性子星を支えるだけの圧力を稼げない（ハイペロンパズルを解決できない）という結論が導かれた場合、状況証拠的に後者の「クォーク物質への相転移シナリオ」の正当性が極めて高く裏付けられることになる。このように、斥力芯の物理を極限まで突き詰めることこそが、相転移の立証に向けた最も確実かつ鋭利なアプローチである。

\subsection{本研究の目的、研究内容、独自性と創造性、どこまで明らかにするか}
\subsubsection{本研究の目的：高密度下での$\Lambda N$相互作用の解明}
本研究の目的は、真空中および高密度媒質中における$\Lambda N$相互作用の短距離成分を実験的に完全解明し、ハイペロンパズルに決着をつけることである。具体的には以下の2つのステップで相互作用の全貌を明らかにする。

\noindent
\textbf{［真空中での$\Lambda N$間の引力獲得メカニズムの理解］}\\
真空中において$\Lambda N$間に引力が生じるミクロな起源として、最も有力視されているのが「$2\pi$交換プロセス」である。これは、$\Lambda$粒子が中間状態で同じくストレンジクォークを1つ含む$\Sigma$粒子へと仮想的に遷移しながら、2つの$\pi$中間子を交換して再び$\Lambda$粒子に戻る二次のプロセスである。このメカニズムが引力の全体像にどの程度寄与しているかを定量的に確定させるためには、反応の頂点における$\Lambda N$-$\Sigma N$間の遷移強度（チャンネル結合力）を実験的に決定しなければならない。本研究では、J-PARC E90実験が推進する「$\Sigma N$カスプ分光」により、この結合力を正確に抽出する。

\noindent
\textbf{［高密度中における$\Lambda N$相互作用の斥力化メカニズムの検証］}\\
真空中では強い引力を生み出す源泉である上記の$2\pi$交換プロセスが、媒質（高密度物質）中に入ると、周囲に存在する核子の密度上昇に伴いパウリ排他律によって強く抑制（パウリブロッキング）される。その結果、中間状態としての$\Sigma$粒子への遷移が阻害されて引力が消滅し、相対的にバリオン内部のクォーク波動関数の重なりに起因する短距離斥力芯が支配的になるという仮説がある。本研究では、通常の原子核の中で特異的に中心密度が高く、標準密度の2倍（$2\rho_0$）に達する$\alpha$粒子（${}^4\mathrm{He}$）に着目する。$\alpha$粒子と$\Lambda$粒子の間の短距離力を「$\alpha\Lambda$フェムトスコピー」によって抽出し、高密度環境特有の斥力の効果を検証する。さらに、研究協力者である三輪らが推進するSPring-8でのHYPS実験やJ-PARC E86実験から得られる中高エネルギー領域の$\Lambda p$散乱データを統合することで真空中の短距離力の絶対基準を確立し、理論分担者（神野・神谷）によるカイラル有効場の理論（Chiral EFT）計算と詳細比較することで、斥力化仮説の真偽を明らかにする。

\subsubsection{研究内容の詳細と到達目標（どこまで明らかにするか）}

\noindent
\textbf{1. $\Sigma N$カスプ分光（E90実験）による$\Lambda N$-$\Sigma N$結合力の決定（分担者：市川、谷田）}\\
およそ50年前、$K^- d \to \Lambda p \pi^-$ 反応において終状態の$\Lambda p$不変質量分布を測定した際、$\Sigma N$生成閾値（$\sim 2130\,\text{MeV}$）の直下に極めて鋭いピーク構造が観測された。これが今日「$\Sigma N$カスプ」と呼ばれる構造である。このカスプ構造は、初期反応で生成された$\Sigma N$対が終状態相互作用によって中間状態で $\Sigma N \to \Lambda N$ へと強くチャンネル転換することに起因して生成される。

理論的には、$\Sigma N \to \Lambda N$ の遷移振幅 $T$ は散乱長を用いて以下のように記述される。
\begin{equation}
    T(\Sigma N \to \Lambda N) \propto \frac{\operatorname{Im}(A_\Sigma)}{1 - i k_\Sigma A_\Sigma}
\end{equation}
ここで、$A_\Sigma$ は低エネルギー$\Sigma N$散乱を特徴づける複素散乱長であり、$A_\Sigma = a + i b$ と表される。実部 $a$ は $\Sigma N$ 間の弾性散乱による引力・斥力のポテンシャル強度を反映し、虚部 $b$ は非弾性過程である $\Sigma N \to \Lambda N$ 遷移の強さ、すなわち本研究が標的とする\textbf{「$\Lambda N$と$\Sigma N$間の結合力」}を直接的に表す物理量である。$k_\Sigma$ は $\Sigma N$ 系における相対運動量である。

実験で観測されるスペクトル強度（反応率）$R$ は遷移振幅の絶対値の2乗 $|T|^2$ に比例するため、閾値近傍における振る舞いは以下のように展開される。
\begin{equation}
    R \propto \frac{1}{|(1 + i k_\Sigma b) - i k_\Sigma a|^2}
\end{equation}
$\Sigma N$閾値よりも上のエネルギー領域（$E = M(\Lambda p) > M_\Sigma + M_N$）では、$k_\Sigma \sim \sqrt{2\mu E}$（$\mu$ は換算質量）となり、反応率は以下の形をとる。
\begin{equation}
    R \sim \frac{4\pi b}{(1 + k_\Sigma b)^2 + (k_\Sigma a)^2}
\end{equation}
一方で、閾値よりも下のエネルギー領域では中間状態の$\Sigma N$は仮想的な粒子となり、運動量は解析接続によって純虚数 $k_\Sigma \to i|k_\Sigma|$ へと置き換わる。これによりスペクトル強度は以下のように記述される。
\begin{equation}
    R \sim \frac{4\pi b}{(1 - |k_\Sigma| a)^2 + (|k_\Sigma| b)^2}
\end{equation}

散乱長パラメータ $a, b$ に様々な値を代入して計算されたカスプスペクトルの理論的な振る舞いからも明らかなように、閾値前後で現れる特徴的かつ非対称なスペクトル形状は散乱長 $A_\Sigma = a+ib$ の値に対して極めて高い感度を有している。したがって、実験的にカスプ形状を高精度で測定し理論式を逆解きすることで、散乱長の実部と虚部を同時に分離して導出することが可能となる。

本研究で推進するJ-PARC E90実験では、高運動量分解能を誇るK1.8ビームラインスペクトロメータおよび新設されたS-2Sスペクトロメータを駆使し、入射 $K^-$ と散乱 $\pi^-$ の運動量を極限精度で測定する。これによりミッシングマス法を用いて$\Sigma N$カスプ質量スペクトルを過去最高の分解能 $\sigma \sim 0.4\,\text{MeV}$ で捉える。さらに、大立体角を持つ三次元飛跡検出器HypTPCを連動させ、崩壊した$\Lambda p$対を同時検出することで準自由過程に起因するバックグラウンドを徹底的に排除する。従来の泡箱実験では統計量と分解能が圧倒的に不足しており散乱長決定には至らなかったが、E90実験は大強度ビームと最新鋭検出器の融合によりこの限界を打破し、散乱長を $0.3\,\text{fm}$ の世界最高精度で確定させる。本実験は分担者の市川および谷田が共同で実験責任者として提案・主導する中核プロジェクトである。

\vspace*{1zw}
\noindent
\textbf{2. $\alpha\Lambda$フェムトスコピーによる高密度媒質中における$\Lambda N$短距離力の解明（分担者：市川・佐甲・大浦）}

\noindent
\textbf{［$\alpha\Lambda$相互作用に関する現状の理解と限界］}\\
$\alpha$粒子と$\Lambda$粒子の間のポテンシャルに関する知見は、長年 ${}^5_\Lambda\mathrm{He}$ ハイパー核の基底状態束縛エネルギー（$B_\Lambda \simeq 3.1\,\text{MeV}$）を再現するという条件によってのみ制限されてきた。しかし、束縛エネルギーという空間積分値のみを制約条件とした場合、短距離領域におけるポテンシャルの詳細な形状には甚大な不定性が残る。複数の代表的な理論モデルはいずれも $B_\Lambda = 3.1\,\text{MeV}$ を完璧に再現するが、近距離における振る舞いは全く異なる。特に一部のモデルは短距離領域において極めて強い斥力芯の存在を示唆する一方で、中心部まで強い引力を維持するモデルも存在している。前述の通り、ハイペロンパズルを$\Lambda N$相互作用の斥力化によって解決するためには、強い短距離斥力芯の実在を実験的に証明しなければならない。

\noindent
\textbf{［極小ソースを用いたフェムトスコピーによる短距離力の抽出］}\\
近年、高エネルギー重イオン衝突等で放出されるハドロン対の相対運動量相関を測定し、二体系の相互作用を調べる「フェムトスコピー」が急速に発展している。二粒子間の相対運動量 $k^*$ に対する実データ（Real event）と混合事象（Mixed event）の比を取ることで相関関数 $C(k^*) = \text{Real}(k^*) / \text{Mixed}(k^*)$ を得る。この関数は粒子放出源の空間サイズ $R$ と二体波動関数に鋭敏に依存する。

分担者の神野・神谷は、これら多様なポテンシャルを用いた $\alpha\Lambda$ 相関関数の精密な理論計算を行った（PRC 110, 014001 (2024)）。その結果、放出源サイズが $R = 3\,\text{fm}$ や $5\,\text{fm}$ と大きい場合、長距離の波動関数が寄与を支配するため全モデルの相関関数は縮退し区別がつかないことが示された。しかし、$R = 1\,\text{fm}$ という極小の放出源サイズを実現した場合、波動関数が短距離ポテンシャルの違いを鋭く拾い上げ、モデル間で相関関数に決定的な差異が生じることが判明した。したがって、極小ソース環境下で相関関数を測定できれば、高密度物質である$\alpha$粒子内部での短距離斥力芯を直接抽出できる。

\noindent
\textbf{［実験計画：J-PARC E88およびRHIC-STARでの測定］}\\
本研究では、J-PARC E88実験（$pA$衝突）のバイプロダクト測定、および米国ブルックヘブン国立研究所のRHIC-STAR実験（エネルギー走査プログラム BES; $AA$衝突）において $\alpha\Lambda$ フェムトスコピーを実行する。$\alpha$粒子のような複合核子は、高エネルギー衝突（$\sqrt{s_{NN}} \gtrsim 200\,\text{GeV}$）では生成率が極めて低いが、$\sqrt{s_{NN}} \approx 10\,\text{GeV}$ 前後の中高エネルギー領域においては生成断面積が100〜1000倍へと爆発的に増大する。

J-PARC E88実験は標的原子核内でのベクター中間子の質量変化を追う実験であるが、大立体角で荷電粒子を捉えるため主目的外の $\alpha\Lambda$ 対も大量に生成される。しかし既存のハードウェアトリガー方式では特定の中間子事象のみを選択するため記録が難しい。そこで本研究では、近年素粒子原子核実験で革命をもたらしている「連続読み出し（ストリーミング）DAQシステム」を独自開発・導入する。これにより全事象をソフトウェアトリガーで柔軟に取得可能とし、極小ソースサイズ $R \sim 1\,\text{fm}$ が保証される $pA$ 反応からの貴重なデータを確実に捉える。また、市川・大浦が参画するSTAR-BES実験（2026年2月データ取得完了）においても周辺衝突事象をセレクトした解析を進めるが、系が非対称になり実効核子数が減る懸念があるため、あくまで本命はJ-PARC E88実験である。

\vspace*{1zw}
\noindent
\textbf{3. 実験データを基にした理論研究（分担者：神野・神谷）}\\
実験から得られた「$\Sigma N$-$\Lambda N$結合力」および「$\alpha\Lambda$短距離斥力」の定量的データを、分担者が主導するカイラル有効場の理論（Chiral EFT）の枠組みへと統合する。低エネルギー定数を最新データで厳密に制限し、標準密度の数倍に至る高密度領域まで適用可能な「密度依存性を持つ包括的バリオン間相互作用モデル」を構築する。これにより中性子星の最大質量予測を行い、パズル解決の可否に最終的な理論的結論を下す。

\subsection{本研究の学術的独自性と創造性}
従来、核力やハドロン間相互作用の研究はビームと標的を用いた散乱実験が主役であった。散乱振幅 $f$ は相対運動量 $k$ と散乱位相差 $\delta$ を用いて $f = 1 / (k\cot\delta - i k)$ と記述され、低エネルギーでは有効レンジ展開 $k\cot\delta \approx 1/a - r_e k^2 / 2$ を介して断面積測定から散乱長 $a$ を導出する。しかし、短寿命なハイペロンを標的とすることは不可能であり、散乱実験の適用範囲は極めて限定的であった。

これに対し、2010年代以降のLHC-ALICE実験等の重イオン衝突によるフェムトスコピーは散乱実験の限界を突破したが、ハドロン対のスピン状態ごとの寄与を分離できないという本質的な原理的制約を抱えている。本研究の第1の独自性と創造性は、半世紀前に理論的枠組みが示されながら技術的限界で放置されていた「カスプ分光」を最新技術で蘇らせる点にある。標的としてスピン $S=1$ を持つ重水素を用いることで、フェムトスコピーには不可能な\textbf{「特定のスピン状態を選択した散乱長の決定」}を実現する。

第2の創造性は、フェムトスコピーの適用領域を拡張し、従来のハドロン同士の相関にとどまらず\textbf{「原子核（$\alpha$粒子）とハドロンの相関」}に踏み込む点にある。中心密度 $2\rho_0$ を持つ実体的な原子核プローブと極小ソースサイズの相乗効果により、長距離引力の寄与を相対的に抑え込み、超高密度下での短距離斥力芯のシグナルのみを劇的に「増幅」させて抽出する。この斬新な手法によって初めて物性の根源たる媒質効果にメスを入れる。

\subsection{関連する国内外の研究動向と本研究の位置づけ}
真空中における $\Lambda N$ 相互作用の長距離成分については、欧州CERNのALICE実験や米国のSTAR実験によるフェムトスコピー研究が世界的に隆盛を極めている。また、中高エネルギー領域（短距離成分）の直接検証としては、研究協力者の三輪らによるSPring-8での高精度 $\Lambda p$ 散乱実験（HYPS）やJ-PARC E86実験が鋭意進行中であり、今後数年で真空中の相互作用に関する包括的なデータが出揃う見込みである。

しかし、天体物理学が直面するハイペロンパズルの解決には「媒質中での相互作用変化」の解明が不可欠である。高エネルギー衝突実験はソースサイズが大きく（$R > 3\,\text{fm}$）、かつ高温からのプラズマ冷却過程であるため、中性子星が要求する「低温・高密度」の物理に直接アクセスできない。本研究は、J-PARCの強みである静止標的実験の特性を活かして放出源を $1\,\text{fm}$ 以下に絞り込み、真空中の引力源泉（結合力）と高密度媒質中の斥力芯を同時に決定する。世界的な研究の進展と完璧に相補的な関係を築きつつ、斥力化シナリオによるパズル解決の決定打を放つ\textbf{唯一無二の国際的ポジション}を確立している。

\subsection{本研究の目的を達成するための準備状況}
本研究の主要実験拠点であるJ-PARCハドロン実験施設K1.8ビームラインは安定稼働を続けており、世界最高強度のストレンジネスビームの利用が保証されている。
\begin{itemize}
    \item \textbf{J-PARC E90実験}：2022年にStage-1承認を取得し、2025年5月に詳細な技術設計報告書（TDR）を提出済みである。現在Stage-2承認に向け、S-2Sスペクトロメータの分解能向上を目指した磁場マップの精密校正や検出器内部の真空化等のハード改良を完了しつつあり、2028年中の物理データ取得に向けて準備は最終段階にある。
    \item \textbf{国際共同研究体制と散乱実験基盤}：理論分担者の神谷はALICE実験Run2データの相関関数解析において顕著な成果を上げており、Run3以降も強固な連携体制を維持する。また真空中の散乱データを補強するため、HYPS実験に向けた革新的な液体水素TPCの開発も並行して進行中である。
    \item \textbf{J-PARC E88およびSTAR-BES実験}：E88実験は既にStage-1承認を得ており、本研究が主導してストリーミングDAQを導入するための読み出し回路およびトリガーファームウェアの開発に着手している。2026年中にバイプロダクト測定に特化した追加実験提案書を提出し、2029年のビームタイムで測定を実施する。一方、STAR-BES実験は2026年2月に全衝突データの取得を完了した。本研究経費によりフェムトスコピー解析に専念する優秀なポスドク（PD）を雇用し、E88実験群と緊密に連携しながら即座に物理解析をスタートさせる万全の体制が整っている。
\end{itemize}

%end 研究目的と研究計画	 ＝＝＝＝＝＝＝＝＝＝＝＝＝＝＝＝＝＝＝＝

\input{pieces/p01_purpose_plan_01}

%#Split: 02_rights  
%#PieceName: p02_rights
\input{pieces/p02_rights_00}
\section{２　人権の保護及び法令等の遵守への対応}
%　　　　＜＜最大　1ページ＞＞

% s09_rights
%begin 人権の保護及び法令等の遵守への対応 ＝＝＝＝＝＝＝＝＝＝＝＝＝＝＝＝＝＝＝＝

%end 人権の保護及び法令等の遵守への対応 ＝＝＝＝＝＝＝＝＝＝＝＝＝＝＝＝＝＝＝＝

\input{pieces/p02_rights_01}

%#Split: 03_final_year  
%#PieceName: p03_final_year
\input{pieces/p03_final_year_00}
\section{３　研究計画最終年度前年度応募を行う場合の記述事項}
%　　　　＜＜最大　1ページ＞＞

%s04_prep_finalyear
% 2020-08-16: Taku: Adjusted the horizontal position of the tabular.
%begin 最終年度の研究課題　＝＝＝＝＝＝＝＝＝＝＝＝＝＝＝＝＝＝＝＝
\newcommand{\最終年度研究種目名}{基盤研究（）}
\newcommand{\最終年度研究課題番号}{?????}
\newcommand{\最終年度研究課題名}{}
\newcommand{\最終年度研究期間}{令和？？年度〜令和\一年目 年度}
%end 最終年度の研究課題　＝＝＝＝＝＝＝＝＝＝＝＝＝＝＝＝＝＝＝＝
\input{pieces/p03_final_year_01}

\noindent
\textbf{当初研究計画及び研究成果}\\
%begin 研究計画最終年度の応募の計画と成果 ＝＝＝＝＝＝＝＝＝＝＝＝＝＝＝＝＝＝＝＝

%end 研究計画最終年度の応募の計画と成果 ＝＝＝＝＝＝＝＝＝＝＝＝＝＝＝＝＝＝＝＝


\noindent
\textbf{前年度応募する理由}\\
%begin 研究計画最終年度の応募の理由 ＝＝＝＝＝＝＝＝＝＝＝＝＝＝＝＝＝＝＝＝

%end 研究計画最終年度の応募の理由 ＝＝＝＝＝＝＝＝＝＝＝＝＝＝＝＝＝＝＝＝

\input{pieces/p03_final_year_02}

%#Split: members_info % keep

% 研究代表者の調書 (not 長所)　＝＝＝＝＝＝＝＝＝＝＝＝＝＝＝
\KLInput{daihyo_info}

% 研究分担者の調書　＝＝＝＝＝＝＝＝＝＝＝＝＝＝＝
\KLInput{buntansha_info}

% 研究分担者を追加する場合は、buntansha_info.tex を複製して別の名前にし、
% この下に追加してください。
% \input ではなく、\KLInput を用いると、読み込んだファイルの中で定義した
% マクロはそのファイルの中でのみ有効となり、他のファイルでの定義と干渉しません。
%\KLInput{taro_info}
%\KLInput{hanako_info}
%\KLInput{jiro_info}
%\KLInput{saburo_info}


%#Split: 99_tail
\input{pieces/hook9} % pieces
\end{document}
